\chapter{2008年山东卷（文）}

\section{单选题}

\begin{problem}
满足 \(M \subseteq \{a_1, a_2, a_3, a_4\}\)，且 \(M \cap \{a_1, a_2, a_3\} = \{a_1, a_2\}\) 的集合 \(M\) 的个数是（\quad）

\choices
    {\(1\)}
    {\(2\)}
    {\(3\)}
    {\(4\)}
\end{problem}

\begin{answer}
B
\end{answer}

\begin{solution}
集合 \(M\) 必须含有 \(a_1\)、\(a_2\)，且不能含有 \(a_3\). 元素 \(a_4\) 可以取，也可以不取，因此
\[
M=\{a_1,a_2\}\quad\text{或}\quad M=\{a_1,a_2,a_4\}
\]
所以满足条件的集合共有 \(2\) 个，选 B.
\end{solution}

\begin{problem}
设 \(z\) 的共轭复数是 \(\overline{z}\)，若 \(z + \overline{z} = 4\)，\(z\overline{z} = 8\)，则 \(\frac{\overline{z}}{z}\) 等于（\quad）

\choices
    {\(\i\)}
    {\(-\i\)}
    {\(\pm 1\)}
    {\(\pm\i\)}
\end{problem}

\begin{answer}
D
\end{answer}

\begin{solution}
设 \(z=x+y\i\)，则
\[
z+\overline{z}=2x=4
\]
，所以 \(x=2\). 又
\[
z\overline{z}=x^2+y^2=8
\]
给出 \(y^2=4\)，即 \(y=\pm2\). 于是 \(z=2\pm2\i\)，从而
\[
\frac{\overline z}{z}
 =\frac{2\mp2\i}{2\pm2\i}
 =\mp\i
\]
故 \(\frac{\overline z}{z}=\pm\i\)，选 D.
\end{solution}

\begin{problem}
函数 \(y = \ln \cos x\)（\(-\frac{\pi}{2} < x < \frac{\pi}{2}\)）的图象是（\quad）

\choices
    {\choicebitmap[width=0.15\paperwidth]{img/2008/shandong_liberal/q3_optA_nolabel.png}}
    {\choicebitmap[width=0.15\paperwidth]{img/2008/shandong_liberal/q3_optB_nolabel.png}}
    {\choicebitmap[width=0.15\paperwidth]{img/2008/shandong_liberal/q3_optC_nolabel.png}}
    {\choicebitmap[width=0.15\paperwidth]{img/2008/shandong_liberal/q3_optD_nolabel.png}}
\end{problem}

\begin{answer}
A
\end{answer}

\begin{solution}
在给定定义域内，
\[
f(-x)=\ln\cos(-x)=\ln\cos x=f(x)
\]
，所以图象关于 \(y\) 轴对称. 又因为 \(0<\cos x\leqslant1\)，所以 \(f(x)\leqslant0\)，且 \(f(0)=0\). 当 \(x\to\pm\frac{\pi}{2}\) 时，\(\cos x\to0^+\)，故 \(f(x)\to-\infty\). 因此图象在原点取得最大值，向两端均下降，符合选项 A.
\end{solution}

\begin{problem}
给出命题：若函数 \( y = f(x) \) 是幂函数，则函数 \( y = f(x) \) 的图象不过第四象限. 在它的逆命题，否命题，逆否命题三个命题中，真命题的个数是（\quad）

\choices
    {\(3\)}
    {\(2\)}
    {\(1\)}
    {\(0\)}
\end{problem}

\begin{answer}
C
\end{answer}

\begin{solution}
设原命题为 \(P\Rightarrow Q\)，其中 \(P\) 表示“\(f\) 是幂函数”，\(Q\) 表示“图象不过第四象限”. 幂函数在其定义域内的图象不会出现 \(x>0\)，\(y<0\) 的点，因此原命题为真，故其逆否命题也为真.

逆命题 \(Q\Rightarrow P\) 不成立，例如函数 \(y=x^2+1\) 的图象不过第四象限，但它不是幂函数. 否命题 \(\lnot P\Rightarrow\lnot Q\) 与逆命题等价，也为假. 因此三个命题中只有一个真命题，选 C.
\end{solution}

\begin{problem}
设函数 \( f(x) = \begin{cases}
1 - x^2, & x \leqslant 1,\\
x^2 + x - 2, & x > 1,
\end{cases} \) 则 \( f\left(\frac{1}{f(2)}\right) \) 的值为（\quad）

\choices
    {\(\frac{15}{16}\)}
    {\( -\frac{27}{16} \)}
    {\( \frac{5}{9} \)}
    {\(18\)}
\end{problem}

\begin{answer}
A
\end{answer}

\begin{solution}
因 \(2>1\)，先用第二段定义计算，得 \(f(2)=2^2+2-2=4\)，所以 \(\frac{1}{f(2)}=\frac14\). 由于 \(\frac14\leqslant1\)，再用第一段定义，得 \(f\left(\frac14\right)=1-\left(\frac14\right)^2=\frac{15}{16}\).
故选 A.
\end{solution}

\begin{problem}
如图是一个几何体的三视图，根据图中数据，可得该几何体的表面积是（\quad）

\bitmapfigure[max width=0.5\linewidth]{img/2008/shandong_liberal/q6_fig1.png}

\choices
    {\(9\pi\)}
    {\(10\pi\)}
    {\(11\pi\)}
    {\(12\pi\)}
\end{problem}

\begin{answer}
D
\end{answer}

\begin{solution}
由三视图可知，该几何体由半径为 \(1\)、高为 \(3\) 的圆柱和半径为 \(1\) 的球组成，球与圆柱顶面相切，因此圆柱的两个底面均属于外表面. 故表面积为
\[
S_{\text{球}}+S_{\text{圆柱侧}}+2S_{\text{圆柱底}}
 =4\pi+2\pi\cdot1\cdot3+2\pi\cdot1^2
 =12\pi
\] 故选 D.
\end{solution}

\begin{problem}
不等式 \(\frac{x + 5}{(x - 1)^2} \geqslant 2\) 的解集（\quad）

\choices
    {\(\left[-3, \frac{1}{2}\right]\)}
    {\(\left[-\frac{1}{2}, 3\right]\)}
    {\(\left[\frac{1}{2}, 1\right) \cup (1, 3]\)}
    {\(\left[-\frac{1}{2}, 1\right) \cup (1, 3]\)}
\end{problem}

\begin{answer}
D
\end{answer}

\begin{solution}
因为 \((x-1)^2>0\)（\(x\neq1\)），不等式等价于 \(x+5\geqslant2(x-1)^2\)，即 \(2x^2-5x-3\leqslant0\). 因式分解得 \((2x+1)(x-3)\leqslant0\)，所以 \(-\frac12\leqslant x\leqslant3\). 再去掉原式无意义的 \(x=1\)，解集为 \(\left[-\frac12,1\right)\cup(1,3]\).
故选 D.
\end{solution}

\begin{problem}
已知 \(a\)，\(b\)，\(c\) 为 \(\triangle ABC\) 的三个内角 \(A\)，\(B\)，\(C\) 的对边，向量 \(\bs{m} = (\sqrt{3}, -1)\)，\(\bs{n} = (\cos A, \sin A)\). 若 \(\bs{m} \perp \bs{n}\)，且 \(a \cos B + b \cos A = c \sin C\)，则角 \(A\)，\(B\) 的大小分别为（\quad）

\choices
    {\(\frac{\pi}{6}\)，\(\frac{\pi}{3}\)}
    {\(\frac{2\pi}{3}\)，\(\frac{\pi}{6}\)}
    {\(\frac{\pi}{3}\)，\(\frac{\pi}{6}\)}
    {\(\frac{\pi}{3}\)，\(\frac{\pi}{3}\)}
\end{problem}

\begin{answer}
C
\end{answer}

\begin{solution}
由 \(\bs{m}\perp\bs{n}\)，得 \(\bs{m}\cdot\bs{n}=\sqrt3\cos A-\sin A=0\). 由于 \(0<A<\pi\)，所以 \(\tan A=\sqrt3\)，从而 \(A=\frac\pi3\).

由正弦定理 \(a=2R\sin A\)，\(b=2R\sin B\)，\(c=2R\sin C\)，原等式化为 \(2R(\sin A\cos B+\sin B\cos A)=2R\sin^2C\). 左边为 \(2R\sin(A+B)=2R\sin C\)，故 \(\sin C=\sin^2C\). 三角形中 \(0<C<\pi\)，所以 \(\sin C>0\)，只能有 \(\sin C=1\)，即 \(C=\frac\pi2\). 于是 \(B=\pi-A-C=\frac\pi6\).
故 \(A=\frac\pi3\)，\(B=\frac\pi6\)，选 C.
\end{solution}

\begin{problem}
从某项综合能力测试中抽取 \(100\) 人的成绩，统计如表，则这 \(100\) 人成绩的标准差为（\quad）

\begin{center}
\begin{tabular}{c|ccccc}
分数 & \(5\) & \(4\) & \(3\) & \(2\) & \(1\) \\
\hline
人数 & \(20\) & \(10\) & \(30\) & \(30\) & \(10\)
\end{tabular}
\end{center}

\choices
    {\(\sqrt{3}\)}
    {\(\frac{2\sqrt{10}}{5}\)}
    {\(3\)}
    {\(\frac{8}{5}\)}
\end{problem}

\begin{answer}
B
\end{answer}

\begin{solution}
设成绩的平均数为 \(\overline{x}\)，则
\[
\overline{x}=\frac{5\cdot20+4\cdot10+3\cdot30+2\cdot30+1\cdot10}{100}=3
\]
方差为
\[
\frac{20(5-3)^2+10(4-3)^2+30(3-3)^2+30(2-3)^2+10(1-3)^2}{100}
 =\frac{160}{100}=\frac85
\]
所以标准差为
\[
\sqrt{\frac85}=\frac{2\sqrt{10}}5
\]选 B.
\end{solution}

\begin{problem}
已知 \(\cos \left( \alpha - \frac{\pi}{6} \right) + \sin \alpha = \frac{4}{5} \sqrt{3}\)，则 \(\sin \left( \alpha + \frac{7\pi}{6} \right)\) 的值是（\quad）

\choices
    {\(-\frac{2\sqrt{3}}{5}\)}
    {\(\frac{2\sqrt{3}}{5}\)}
    {\(-\frac{4}{5}\)}
    {\(\frac{4}{5}\)}
\end{problem}

\begin{answer}
C
\end{answer}

\begin{solution}
展开已知等式，得
\[
\frac{\sqrt3}{2}\cos\alpha+\frac12\sin\alpha+\sin\alpha
 =\frac{4\sqrt3}{5}
\]
两边乘以 \(\frac2{\sqrt3}\)，得到
\[
\cos\alpha+\sqrt3\sin\alpha=\frac85
\]
另一方面，
\[
\sin\left(\alpha+\frac{7\pi}{6}\right)
 =-\frac12\cos\alpha-\frac{\sqrt3}{2}\sin\alpha
 =-\frac12\left(\cos\alpha+\sqrt3\sin\alpha\right)
 =-\frac45
\]
故选 C.
\end{solution}

\begin{problem}
若圆 \(C\) 的半径为 \(1\)，圆心在第一象限，且与直线 \(4x - 3y = 0\) 和 \(x\) 轴相切，则该圆的标准方程是（\quad）

\choices
    {\((x - 3)^{2} + \left(y - \frac{7}{3}\right)^{2} = 1\)}
    {\((x - 2)^{2} + (y - 1)^{2} = 1\)}
    {\((x - 1)^{2} + (y - 3)^{2} = 1\)}
    {\(\left(x - \frac{3}{2}\right)^2 + (y - 1)^2 = 1\)}
\end{problem}

\begin{answer}
B
\end{answer}

\begin{solution}
设圆心为 \((h,k)\). 圆与 \(x\) 轴相切且半径为 \(1\)，圆心在第一象限，所以 \(k=1\). 又圆心到直线 \(4x-3y=0\) 的距离为 \(1\)，因此 \(\frac{|4h-3|}{5}=1\)，即 \(|4h-3|=5\). 解得 \(h=2\) 或 \(h=-\frac12\). 因圆心在第一象限，取 \(h=2\)，故圆的标准方程为
\[
(x-2)^2+(y-1)^2=1
\]
选 B.
\end{solution}

\begin{problem}
已知函数 \( f(x)=\log_{a}(2^{x}+b-1) \)（\(a>0\)，\(a\neq1\)）的图象如图所示，则 \(a\)，\(b\) 满足的关系是（\quad）

\bitmapfigure[max width=6.0cm]{img/2008/shandong_liberal/q12_fig1.png}

\choices
    {\(0 < a^{-1} < b < 1\)}
    {\(0 < b < a^{-1} < 1\)}
    {\(0 < b^{-1} < a < -1\)}
    {\(0 < a^{-1} < b^{-1} < 1\)}
\end{problem}

\begin{answer}
A
\end{answer}

\begin{solution}
由图象单调递增，知 \(a>1\). 图象在 \(y\) 轴上的截距位于 \(-1\) 与 \(0\) 之间，所以
\[
-1< f(0)=\log_a b<0
\]
由于 \(a>1\) 时对数函数单调递增，且 \(\log_a(a^{-1})=-1\)、\(\log_a1=0\)，故
\[
a^{-1}<b<1
\]
结合 \(b>0\)，得到 \(0<a^{-1}<b<1\)，选 A.
\end{solution}

\section{填空题}

\begin{problem}
已知圆 \(C: x^{2} + y^{2} - 6x - 4y + 8 = 0\). 以圆 \(C\) 与坐标轴的交点分别作为双曲线的一个焦点和顶点，则适合上述条件的双曲线的标准方程为\fillinblank{}.
\end{problem}

\begin{answer}
\(\frac{x^{2}}{4} - \frac{y^{2}}{12} = 1\)
\end{answer}

\begin{solution}
将圆方程配方，得
\[
(x-3)^2+(y-2)^2=5
\]
它与 \(x\) 轴的交点满足 \(y=0\)，故
\[
(x-3)^2+4=5
\]交点为 \((2,0)\)、\((4,0)\). 当 \(x=0\) 时，圆方程化为 \(y^2-4y+8=0\)，其判别式为 \(-16\)，所以圆与 \(y\) 轴没有交点. 对于标准形式的水平双曲线，焦半距满足 \(c>a>0\)，因此只能取 \((2,0)\) 为双曲线的顶点、\((4,0)\) 为同一支的焦点，从而 \(a=2\)，\(c=4\). 由 \(c^2=a^2+b^2\)，得 \(b^2=16-4=12\). 所以标准方程为
\[
\frac{x^2}{4}-\frac{y^2}{12}=1
\]
\end{solution}

\begin{problem}
执行下面的程序框图，若 \(p = 0.8\)，则输出的 \(n =\)\fillinblank{}.

\bitmapfigure[max width=0.72\linewidth]{img/2008/shandong_liberal/q14_fig1.png}
\end{problem}

\begin{answer}
\(4\)
\end{answer}

\begin{solution}
初始时 \(n=1\)，\(S=0\). 逐次判断并更新：
\begin{center}
\begin{tabular}{c|c|c}
\(n\) & \(S\) & \(S<p?\) \\ \hline
\(1\) & \(0\) & \(0<0.8\)，更新为 \(S=\frac12\)，\(n=2\) \\
\(2\) & \(\frac12\) & \(\frac12<0.8\)，更新为 \(S=\frac34\)，\(n=3\) \\
\(3\) & \(\frac34\) & \(\frac34<0.8\)，更新为 \(S=\frac34+\frac18=\frac78\)，\(n=4\) \\
\(4\) & \(\frac78\) & \(\frac78=0.875>0.8\)，停止
\end{tabular}
\end{center}
因此输出 \(n=4\).
\end{solution}

\begin{problem}
已知 \( f(3^{x}) = 4x \log_{2} 3 + 233 \)，则 \( f(2) + f(4) + f(8) + \cdots + f(2^{8}) \) 的值等于\fillinblank{}.
\end{problem}

\begin{answer}
\(2008\)
\end{answer}

\begin{solution}
令 \(t=3^x\)，则 \(x=\log_3t\). 由换底公式，\(\log_2 3\cdot\log_3t=\log_2t\)，所以 \(f(t)=4\log_2t+233\).
当 \(t=2^k\) 时，\(f(2^k)=4k+233\). 因此
\[
f(2)+f(4)+\cdots+f(2^8)
=\sum_{k=1}^{8}(4k+233)
=4\cdot\frac{8\cdot9}{2}+8\cdot233
 =2008
\]
\end{solution}

\begin{problem}
设 \(x\)，\(y\) 满足约束条件 \(\begin{cases}
x - y + 2 \geqslant 0,\\
5x - y - 10 \leqslant 0,\\
x \geqslant 0,\\
y \geqslant 0.
\end{cases}\) 则 \(z = 2x + y\) 的最大值为\fillinblank{}.
\end{problem}

\begin{answer}
\(11\)
\end{answer}

\begin{solution}
约束条件给出可行域的边界 \(y=x+2\)，\(y=5x-10\)，\(x=0\)，\(y=0\). 其顶点为 \((0,0)\)、\((0,2)\)、两直线交点 \((3,5)\) 以及 \((2,0)\). 在这些顶点处计算目标函数 \(z=2x+y\)，所得值依次为 \(0\)，\(2\)，\(11\)，\(4\).
线性目标函数在多边形顶点处取得最大值，所以最大值为 \(11\).
\end{solution}

\section{解答题}

\begin{problem}
已知函数 \( f(x) = \sqrt{3}\sin (\omega x + \varphi) - \cos (\omega x + \varphi) \)（\(0 < \varphi < \pi\)，\(\omega > 0\)）为偶函数，且函数 \( y = f(x) \) 图象的两相邻对称轴间的距离为 \( \frac{\pi}{2} \).
\begin{enumerate}
\item 求 \( f\left(\frac{\pi}{8}\right) \) 的值；
\item 将函数 \( y = f(x) \) 的图象向右平移 \( \frac{\pi}{6} \) 个单位后，得到函数 \( y = g(x) \) 的图象，求 \( g(x) \) 的单调递减区间.
\end{enumerate}
\end{problem}

\begin{solution}
\begin{enumerate}
\item 先合并同角三角函数，得 \(f(x)=2\sin\left(\omega x+\varphi-\frac{\pi}{6}\right)\). 令 \(\delta=\varphi-\frac\pi6\). 因 \(f\) 为偶函数，对任意 \(x\) 有 \(0=f(x)-f(-x)=4\cos\delta\sin(\omega x)\). 由于 \(\omega>0\)，这对任意 \(x\) 成立只能有 \(\cos\delta=0\). 又 \(-\frac\pi6<\delta<\frac{5\pi}6\)，故 \(\delta=\frac\pi2\)，从而 \(\varphi=\frac{2\pi}3\)，并且 \(f(x)=2\cos(\omega x)\). 函数 \(2\cos(\omega x)\) 的相邻对称轴间距为 \(\frac\pi\omega\)，所以 \(\frac\pi\omega=\frac\pi2\)，即 \(\omega=2\). 于是 \(f\left(\frac\pi8\right)=2\cos\frac\pi4=\sqrt2\).

\item 图象向右平移 \(\frac\pi6\) 后，\(g(x)=f\left(x-\frac\pi6\right)=2\cos\left(2x-\frac\pi3\right)\)，从而 \(g'(x)=-4\sin\left(2x-\frac\pi3\right)\). 因此 \(g\) 递减当且仅当 \(\sin\left(2x-\frac\pi3\right)>0\)，即存在 \(k\in\Z\)，使 \(2k\pi<2x-\frac\pi3<(2k+1)\pi\). 所以 \(g(x)\) 的单调递减区间为 \(\left(k\pi+\frac\pi6,k\pi+\frac{2\pi}3\right)\). 按通常教材将端点并入单调区间的记法，写为 \(\left[k\pi+\frac\pi6,k\pi+\frac{2\pi}3\right]\)，其中 \(k\in\Z\).
\end{enumerate}
\end{solution}

\begin{problem}
现有 \(8\) 名奥运会志愿者，其中志愿者 \(A_{1}\)，\(A_{2}\)，\(A_{3}\) 通晓日语，\(B_{1}\)，\(B_{2}\)，\(B_{3}\) 通晓俄语，\(C_{1}\)，\(C_{2}\) 通晓韩语. 从中选出通晓日语，俄语和韩语的志愿者各 \(1\) 名，组成一个小组.
\begin{enumerate}
\item 求 \(A_{1}\) 被选中的概率；
\item 求 \(B_{1}\) 和 \(C_{1}\) 不全被选中的概率.
\end{enumerate}
\end{problem}

\begin{solution}
\begin{enumerate}
\item 每个小组由一名日语、一名俄语和一名韩语志愿者组成，共有 \(3\cdot3\cdot2=18\) 个等可能的小组. 包含 \(A_1\) 时，俄语和韩语志愿者分别有 \(3\) 种和 \(2\) 种选法，共 \(6\) 个小组，因此 \(P(A_1\text{ 被选中})=\frac6{18}=\frac13\).

\item 同时选中 \(B_1\)、\(C_1\) 时，日语志愿者有 \(3\) 种选法，所以该事件有 \(3\) 个小组，概率为 \(\frac3{18}=\frac16\). 因此 \(B_1\)、\(C_1\) 不全被选中的概率为
\(1-\frac16=\frac56\).
\end{enumerate}
\end{solution}

\begin{problem}
如图，在四棱锥 \(P - ABCD\) 中，平面 \(PAD \perp\) 平面 \(ABCD\)，\(AB \parallel DC\)，\(\triangle PAD\) 是等边三角形，已知 \(BD = 2AD = 8\)，\(AB = 2DC = 4\sqrt{5}\).
\begin{enumerate}
\item 设 \(M\) 是 \(PC\) 上的一点，证明：平面 \(MBD \perp\) 平面 \(PAD\)；
\item 求四棱锥 \(P-ABCD\) 的体积.
\end{enumerate}

\bitmapfigure[max width=0.40\linewidth]{img/2008/shandong_liberal/q19_fig1.png}
\end{problem}

\begin{solution}
\begin{enumerate}
\item 由 \(BD=2AD=8\)，得 \(AD=4\)，\(BD=8\). 又 \(AB=4\sqrt5\)，故
\[
AB^2=(4\sqrt5)^2=80=4^2+8^2=AD^2+BD^2
\]
在三角形 \(ABD\) 中，由勾股定理的逆定理知 \(AD\perp BD\). 直线 \(BD\) 在平面 \(ABCD\) 内，且平面 \(PAD\perp\) 平面 \(ABCD\)，两平面的交线为 \(AD\)，所以平面 \(ABCD\) 内垂直于 \(AD\) 的直线 \(BD\) 垂直于平面 \(PAD\). 由于直线 \(BD\) 在平面 \(MBD\) 内，可得平面 \(MBD\perp\) 平面 \(PAD\).

\item 先求底面面积. 建立平面直角坐标系，取 \(A=(0,0)\)、\(B=(4\sqrt5,0)\)，设 \(D=(u,h)\). 因 \(AB\parallel DC\) 且四边形为凸四边形，可取 \(C=(u+2\sqrt5,h)\). 由 \(AD=4\)、\(BD=8\)，有
\[
\begin{cases}
u^2+h^2=16,\\
(u-4\sqrt5)^2+h^2=64
\end{cases}
\]
两式相减得 \(u=\frac4{\sqrt5}\)，代回得 \(h=\frac8{\sqrt5}\). 所以梯形 \(ABCD\) 的面积为
\[
S_{ABCD}=\frac{4\sqrt5+2\sqrt5}{2}\cdot\frac8{\sqrt5}=24
\]
等边三角形 \(PAD\) 的边长为 \(AD=4\). 设 \(O\) 为 \(AD\) 的中点，则
\[
PO=\sqrt{4^2-2^2}=2\sqrt3
\]
又 \(PO\perp AD\)，且 \(PO\) 在平面 \(PAD\) 内，故由两平面垂直的性质 \(PO\perp\) 平面 \(ABCD\)，它就是四棱锥的高. 因此
\[
V_{P-ABCD}=\frac13S_{ABCD}\cdot PO
 =\frac13\cdot24\cdot2\sqrt3
 =16\sqrt3
\]
\end{enumerate}
\end{solution}

\begin{problem}
将数列 \(\{a_{n}\}\) 中的所有项按每一行比上一行多一项的规则排成如下数表：
\begin{center}
\begin{tabular}{cccc}
\(a_1\) & & & \\
\(a_2\) & \(a_3\) & & \\
\(a_4\) & \(a_5\) & \(a_6\) & \\
\(a_7\) & \(a_8\) & \(a_9\) & \(a_{10}\)
\end{tabular}
\end{center}
记表中的第一列数 \(a_1\)，\(a_2\)，\(a_4\)，\(a_7\)，\(\cdots\) 构成的数列为 \(\{b_n\}\)，\(b_1 = a_1 = 1\). \(S_n\) 为数列 \(\{b_n\}\) 的前 \(n\) 项和，且满足 \(\frac{2b_n}{b_n S_n - S_n^{2}} = 1\)（\(n \geqslant 2\)）.
\begin{enumerate}
\item 证明数列 \(\left\{\frac{1}{S_n}\right\}\) 成等差数列，并求数列 \(\{b_n\}\) 的通项公式；
\item 上表中，若从第三行起，每一行中的数按从左到右的顺序均构成等比数列，且公比为同一个正数. 当 \(a_{81} = -\frac{4}{91}\) 时，求上表中第 \(k\)（\(k \geqslant 3\)）行所有项和的和.
\end{enumerate}
\end{problem}

\begin{solution}
\begin{enumerate}
\item 由 \(b_n=S_n-S_{n-1}\)，原条件可化为
\[
2b_n=b_nS_n-S_n^2=S_n(b_n-S_n)=-S_nS_{n-1}
\]
题设分式有意义，所以 \(S_n\neq0\). 若 \(S_{n-1}=0\)，上式将给出 \(b_n=0\)，进而 \(S_n=S_{n-1}+b_n=0\)，矛盾；故 \(S_nS_{n-1}\neq0\)，可以相除，得到
\[
\frac1{S_n}-\frac1{S_{n-1}}=\frac12
\]
而 \(S_1=b_1=1\)，所以 \(\left\{\frac1{S_n}\right\}\) 是首项为 \(1\)、公差为 \(\frac12\) 的等差数列，
\[
\frac1{S_n}=1+\frac{n-1}{2}=\frac{n+1}{2}
\]
于是 \(S_n=\frac2{n+1}\).

\[
b_n=\begin{cases}
1, & n=1,\\
-\dfrac{2}{n(n+1)}, & n\geqslant 2
\end{cases}
\]
\item 第 \(k\) 行的第一个数是 \(b_k\). 因为前 \(12\) 行共有
\[
1+2+\cdots+12=78
\]
项，而前 \(13\) 行共有 \(91\) 项，所以 \(a_{81}\) 是第 \(13\) 行从左数第 \(3\) 项. 由 \(b_{13}=-\frac1{91}\)，设各行公比为 \(q>0\)，则
\[
a_{81}=b_{13}q^2=-\frac4{91}
\]
从而 \(q^2=4\)，即 \(q=2\). 因此第 \(k\) 行的和为
\(T_k=b_k(1+q+\cdots+q^{k-1}) =-\frac2{k(k+1)}(1+2+\cdots+2^{k-1}) =\frac{2(1-2^k)}{k(k+1)}\)，\((k\geqslant3)\)
\end{enumerate}
\end{solution}

\begin{problem}
设函数 \( f(x) = x^{2}\e^{x - 1} + ax^{3} + bx^{2} \)，已知 \( x = -2 \) 和 \( x = 1 \) 为 \( f(x) \) 的极值点.
\begin{enumerate}
\item 求 \(a\) 和 \(b\) 的值；
\item 讨论 \( f(x) \) 的单调性；
\item 设 \( g(x)=\frac{2}{3}x^{3}-x^{2} \)，试比较 \( f(x) \) 与 \( g(x) \) 的大小.
\end{enumerate}
\end{problem}

\begin{solution}
\begin{enumerate}
\item 求导得
\[
f'(x)=(x^2+2x)\e^{x-1}+3ax^2+2bx
\]
由于 \(-2\)、\(1\) 是极值点，必有 \(f'(-2)=f'(1)=0\). 代入得
\[
\begin{cases}
12a-4b=0,\\
3+3a+2b=0
\end{cases}
\]
解得
\(a=-\frac13\)，\(b=-1\)
\item 代入 \(a\)、\(b\)，并因式分解：
\[
f'(x)=(x^2+2x)\e^{x-1}-x^2-2x
=x(x+2)(\e^{x-1}-1)
\]
所以 \(f'(x)\) 的零点为 \(-2\)、\(0\)、\(1\). 按各因子的符号讨论，得
\begin{center}
\begin{tabular}{c|cccc}
\(x\) & \((-\infty,-2)\) & \((-2,0)\) & \((0,1)\) & \((1,+\infty)\) \\ \hline
\(f'(x)\) & \(-\) & \(+\) & \(-\) & \(+\)
\end{tabular}
\end{center}
故 \(f\) 在 \((-\infty,-2)\)、\((0,1)\) 上单调递减，在 \((-2,0)\)、\((1,+\infty)\) 上单调递增.
\item
\[
f(x)-g(x)
=x^2\e^{x-1}-\frac13x^3-x^2-\left(\frac23x^3-x^2\right)
=x^2(\e^{x-1}-x)
\]
令 \(h(x)=\e^{x-1}-x\)，则 \(h'(x)=\e^{x-1}-1\). 因此 \(h\) 在 \(x=1\) 处取得最小值，且 \(h(1)=0\)，从而 \(h(x)\geqslant0\). 又 \(x^2\geqslant0\)，所以对任意实数 \(x\)，都有 \(f(x)-g(x)\geqslant0\)，即 \(f(x)\geqslant g(x)\)，且等号当且仅当 \(x=0\) 或 \(x=1\).
\end{enumerate}
\end{solution}

\begin{problem}
已知曲线 \(C_1: \frac{|x|}{a} + \frac{|y|}{b} = 1\)（\(a > b > 0\)）所围成的封闭图形的面积为 \(4\sqrt{5}\)，曲线 \(C_1\) 的内切圆半径为 \(\frac{2\sqrt{5}}{3}\). 记 \(C_2\) 为以曲线 \(C_1\) 与坐标轴的交点为顶点的椭圆.
\begin{enumerate}
\item 求椭圆 \(C_{2}\) 的标准方程；
\item 设 \(AB\) 是过椭圆 \(C_{2}\) 中心的任意弦，\(l\) 是线段 \(AB\) 的垂直平分线. \(M\) 是 \(l\) 上异于椭圆中心的点.
\begin{enumerate}
\item 若 \( |MO| = \lambda |OA| \)（\(O\) 为坐标原点），当点 \(A\) 在椭圆 \(C_2\) 上运动时，求点 \(M\) 的轨迹方程；
\item 若 \(M\) 是 \(l\) 与椭圆 \(C_2\) 的交点，求 \(\triangle AMB\) 的面积的最小值.
\end{enumerate}
\end{enumerate}
\end{problem}

\begin{solution}
\begin{enumerate}
\item 曲线 \(C_1\) 与坐标轴的交点为 \((\pm a,0)\)、\((0,\pm b)\)，其围成的图形面积为 \(2ab\). 由面积公式，\(2ab=4\sqrt5\)，所以 \(ab=2\sqrt5\). 该菱形的边长为 \(\sqrt{a^2+b^2}\)，半周长为 \(2\sqrt{a^2+b^2}\)，故内切圆半径满足
\[
\frac{2ab}{2\sqrt{a^2+b^2}}=\frac{2\sqrt5}{3}
\]
代入 \(ab=2\sqrt5\)，得 \(a^2+b^2=9\). 于是 \(a^2\)、\(b^2\) 是方程
\[
t^2-9t+20=0
\]
的两个根，即 \(a^2=5\)，\(b^2=4\). 因此椭圆方程为
\[
\frac{x^2}{5}+\frac{y^2}{4}=1
\]
\item
\begin{enumerate}
\item 设 \(A=(u,v)\)，则 \(B=(-u,-v)\)，且 \(u^2/5+v^2/4=1\). 因为 \(l\) 是 \(AB\) 的垂直平分线，所以 \(OM\perp OA\). 取垂直于 \(OA\) 的方向向量 \((-v,u)\)，由 \(|MO|=\lambda|OA|\) 可设
\[
M=\pm\lambda(-v,u)
\]
若 \(M=(x,y)\)，则 \(x^2=\lambda^2v^2\)、\(y^2=\lambda^2u^2\). 结合椭圆方程，得到
\[
\frac{x^2}{4}+\frac{y^2}{5}
=\lambda^2\left(\frac{v^2}{4}+\frac{u^2}{5}\right)
=\lambda^2
\]
故轨迹方程为 \(\frac{x^2}{4}+\frac{y^2}{5}=\lambda^2\)，其中 \(\lambda>0\).

\item 设 \(M=(u,v)\)，则 \(u^2/5+v^2/4=1\). 弦 \(AB\) 所在直线过原点且垂直于 \(OM\)，其方向向量可取 \((-v,u)\). 令弦端点为 \(\pm t(-v,u)\)，代入椭圆方程得
\[
t^2\left(\frac{v^2}{5}+\frac{u^2}{4}\right)=1
\]
故
\[
AB=\frac{2\sqrt{u^2+v^2}}{\sqrt{v^2/5+u^2/4}}
\]
从而
\[
S_{\triangle AMB}=\frac{u^2+v^2}{\sqrt{v^2/5+u^2/4}}
\]
令 \(s=u^2\)，则 \(0\leqslant s\leqslant5\)，且由椭圆方程 \(v^2=4-\frac45s\). 于是
\[
S_{\triangle AMB}=\frac{2(20+s)}{\sqrt{80+9s}}
\]
令 \(H(s)=S_{\triangle AMB}^{\,2}=\dfrac{4(20+s)^2}{80+9s}\)，则
\[
H'(s)=\frac{4(20+s)(9s-20)}{(80+9s)^2}
\]
因为 \(20+s>0\) 且 \(80+9s>0\)，所以 \(H'(s)\) 的符号由 \(9s-20\) 决定，面积在 \(s=\frac{20}{9}\) 时取得最小值. 代入得
\[
S_{\min}=\frac{2\left(20+\frac{20}{9}\right)}{\sqrt{80+9\cdot\frac{20}{9}}}
=\frac{40}{9}
\]
\end{enumerate}
\end{enumerate}
\end{solution}
